% Options for packages loaded elsewhere
\PassOptionsToPackage{unicode}{hyperref}
\PassOptionsToPackage{hyphens}{url}
\documentclass[
]{article}
\usepackage{xcolor}
\usepackage{amsmath,amssymb}
\setcounter{secnumdepth}{-\maxdimen} % remove section numbering
\usepackage{iftex}
\ifPDFTeX
  \usepackage[T1]{fontenc}
  \usepackage[utf8]{inputenc}
  \usepackage{textcomp} % provide euro and other symbols
\else % if luatex or xetex
  \usepackage{unicode-math} % this also loads fontspec
  \defaultfontfeatures{Scale=MatchLowercase}
  \defaultfontfeatures[\rmfamily]{Ligatures=TeX,Scale=1}
\fi
\usepackage{lmodern}
\ifPDFTeX\else
  % xetex/luatex font selection
\fi
% Use upquote if available, for straight quotes in verbatim environments
\IfFileExists{upquote.sty}{\usepackage{upquote}}{}
\IfFileExists{microtype.sty}{% use microtype if available
  \usepackage[]{microtype}
  \UseMicrotypeSet[protrusion]{basicmath} % disable protrusion for tt fonts
}{}
\makeatletter
\@ifundefined{KOMAClassName}{% if non-KOMA class
  \IfFileExists{parskip.sty}{%
    \usepackage{parskip}
  }{% else
    \setlength{\parindent}{0pt}
    \setlength{\parskip}{6pt plus 2pt minus 1pt}}
}{% if KOMA class
  \KOMAoptions{parskip=half}}
\makeatother
\usepackage{longtable,booktabs,array}
\newcounter{none} % for unnumbered tables
\usepackage{calc} % for calculating minipage widths
% Correct order of tables after \paragraph or \subparagraph
\usepackage{etoolbox}
\makeatletter
\patchcmd\longtable{\par}{\if@noskipsec\mbox{}\fi\par}{}{}
\makeatother
% Allow footnotes in longtable head/foot
\IfFileExists{footnotehyper.sty}{\usepackage{footnotehyper}}{\usepackage{footnote}}
\makesavenoteenv{longtable}
\setlength{\emergencystretch}{3em} % prevent overfull lines
\providecommand{\tightlist}{%
  \setlength{\itemsep}{0pt}\setlength{\parskip}{0pt}}
\usepackage{bookmark}
\IfFileExists{xurl.sty}{\usepackage{xurl}}{} % add URL line breaks if available
\urlstyle{same}
\hypersetup{
  hidelinks,
  pdfcreator={LaTeX via pandoc}}

\author{}
\date{}

\begin{document}

第二章 一元函数微分学

\paragraph{一、导数}\label{ux4e00ux5bfcux6570}

\begin{enumerate}
\def\labelenumi{\arabic{enumi}.}
\item
  \textbf{定义：}
\end{enumerate}

\begin{itemize}
\item
  设函数\(y=f(x)\)在点\(x_0\)的某领域内有定义，当自变量的增量
  \(\Delta x\) 趋于 0 时，函数增量 \(\Delta y\) 与自变量增量
  \(\Delta x\) 之比的极限
\end{itemize}

\[f'(x_0) = \lim_{\Delta x \to 0} \frac{\Delta y}{\Delta x} = \lim_{\Delta x \to 0} \frac{f(x_0 + \Delta x) - f(x_0)}{\Delta x}\]

存在，则称\(y=f(x)\)在\(x_0\)处可导，\textbf{此极限值称为\(f(x)\)在\(x_0\)处的导数}，记作\(f'(x_0)\)，或\(y'|_{x=x_0}\)，\(\frac {dy}{dx}|_{x=x_0}\)。

\begin{itemize}
\item
  令\(x=x_0+\Delta x\)，可得导数的等效定义
\end{itemize}

\[f'(x)=\lim_{x \to x_0} \frac {f(x)-f(x_0)}{x-x_0}\]

\begin{enumerate}
\def\labelenumi{\arabic{enumi}.}
\item
  区间可导
\end{enumerate}

\begin{itemize}
\item
  开区间可导：如果函数 \(f(x)\) 在开区间 \((a, b)\)
  内的\textbf{每一个点}都可导，我们就说 \(f(x)\) 在开区间 \((a, b)\)
  内可导。
\item
  闭区间可导：如果函数 \(f(x)\)
  满足以下\textbf{三个条件}，我们就说它在闭区间 \([a, b]\) 上可导：

  \begin{itemize}
  \item
    在内部开区间 \((a, b)\) 内可导。
  \item
    在左端点 \(a\) 处，\textbf{右导数}存在（即 \(f'_+(a)\) 存在）。
  \item
    在右端点 \(b\) 处，\textbf{左导数}存在（即 \(f'_-(b)\) 存在）。
  \end{itemize}
\end{itemize}

\begin{enumerate}
\def\labelenumi{\arabic{enumi}.}
\item
  函数导数的条件：
\end{enumerate}

\begin{itemize}
\item
  必要条件：函数\(f(x)\)在\(x_0\)处连续。
\item
  充要条件：\(f'_-(x_0)\)和 \( f'_+(x_0)\)存在且相等。
\end{itemize}

\begin{enumerate}
\def\labelenumi{\arabic{enumi}.}
\item
  \textbf{计算}
\end{enumerate}

\begin{itemize}
\item
  基本初等函数的导数
\end{itemize}

{\def\LTcaptype{none} % do not increment counter
\begin{longtable}[]{@{}lll@{}}
\toprule\noalign{}
\textbf{函数类型} & \textbf{原函数 \(f(x)\)} & \textbf{导函数
\(f'(x)\)} \\
\midrule\noalign{}
\endhead
\bottomrule\noalign{}
\endlastfoot
\textbf{常数 / 幂函数} & 略 & \\
\textbf{指数 / 对数} & 略 & \\
\textbf{三角函数} & \(\sin x\) & \(\cos x\) \\
& \(\cos x\) & \(-\sin x\) \\
& \(\tan x\) & \(\sec^2 x\) \\
& \(\cot x\) & \(-\csc^2 x\) \\
& \(\sec x\) & \(\sec x \tan x\) \\
& \(\csc x\) & \(-\csc x \cot x\) \\
\textbf{反三角函数} & \(\arcsin x\) & \(\frac{1}{\sqrt{1-x^2}}\) \\
& \(\arccos x\) & \(-\frac{1}{\sqrt{1-x^2}}\) \\
& \(\arctan x\) & \(\frac{1}{1+x^2}\) \\
\end{longtable}
}

\begin{itemize}
\item
  \textbf{四则运算}

  设 \(u(x)\) 和 \(v(x)\) 均可导：

  \begin{itemize}
  \item
    \textbf{加减法：} \((u \pm v)' = u' \pm v'\)
  \item
    \textbf{乘法：} \((uv)' = u'v + uv'\)
  \item
    \textbf{除法：}
    \(\left(\frac{u}{v}\right)' = \frac{u'v - uv'}{v^2}\) （前提
    \(v \neq 0\)）
  \end{itemize}
\item
  \textbf{复合函数求导（链式法则 Chain Rule）------ 核心中的核心}

  遇到函数套函数，必须\textbf{由外向内，逐层求导，层层相乘}。

  \begin{itemize}
  \item
    如果 \(y = f(u)\)，且 \(u = g(x)\)，那么：

    \[\frac{dy}{dx} = f'(u) \cdot g'(x)\]
  \item
    \textbf{实战忠告：} 很多时候容易漏掉最内层自变量的求导（比如
    \((e^{-x})' = -e^{-x}\)，容易漏掉负号）。
  \end{itemize}
\item
  \textbf{反函数求导}：如果函数 \(y = f(x)\) 在某区间内单调、可导且
  \(f'(x) \neq 0\)，那么它的反函数 \(x = \varphi(y)\)
  在对应区间内也可导，且公式为：

  \[\frac{dx}{dy} = \frac{1}{\frac{dy}{dx}}\]

  \begin{itemize}
  \item
    反函数的二阶导数：
  \end{itemize}

  \[\frac{d^2x}{dy^2} = \frac{d}{dy}\left( \frac{dx}{dy} \right) = \frac{d}{dy}\left( \frac{1}{y'(x)} \right)= \frac{d}{dx}\left( [y'(x)]^{-1} \right) \cdot \frac{dx}{dy}= -[y'(x)]^{-2} \cdot y''(x) \cdot \frac{1}{y'(x)}=\frac{d^2x}{dy^2} = -\frac{y''}{(y')^3}\]
\item
  \textbf{隐函数求导}：设\(y=f(x)\)是由方程 \(F(x,y)=0\)
  所确定的可导函数，则\textbf{将方程 \(F(x,y) = 0\) 的两边同时对 \(x\)
  求导}，可得到一个含有\(y'\)的方程，解出\(y'\)即可。
\item
  \textbf{参数方程求导}：设\(y=f(x)\)是由参数方程\(\begin{cases} x = x(t) \\ y = y(t) \end{cases}\)确定的函数，若\(x(t)\)和\(y(t)\)都可导，则
\end{itemize}

\[\frac{dy}{dx} = \frac{\frac{dy}{dt}}{\frac{dx}{dt}} = \frac{\psi'(t)}{\varphi'(t)}\]

\begin{itemize}
\item
  \textbf{高阶导数}

  \begin{itemize}
  \item
    莱布尼茨公式
  \end{itemize}

  \[(uv)^{(n)} = u^{(n)}v + C_n^1 u^{(n-1)}v' + C_n^2 u^{(n-2)}v'' + \dots + u v^{(n)}\]

  \begin{itemize}
  \item
    常用公式
  \end{itemize}
\end{itemize}

{\def\LTcaptype{none} % do not increment counter
\begin{longtable}[]{@{}lll@{}}
\toprule\noalign{}
\textbf{函数 f(x)} & \textbf{n 阶导数 f(n)(x)} & \textbf{记忆特征} \\
\midrule\noalign{}
\endhead
\bottomrule\noalign{}
\endlastfoot
\textbf{\(e^{kx}\)} & \(k^n e^{kx}\) & 每次求导掉下一个 \(k\) \\
\textbf{\(\sin(kx)\)} & \(k^n \sin(kx + \frac{n\pi}{2})\) &
每次求导提一个 \(k\)，相位加 \(\frac{\pi}{2}\) \\
\textbf{\(\cos(kx)\)} & \(k^n \cos(kx + \frac{n\pi}{2})\) & 同上 \\
\textbf{\(\ln(1+x)\)} & \((-1)^{n-1} \frac{(n-1)!}{(1+x)^n}\) &
符号交替，分子是 \((n-1)!\) \\
\textbf{\(\frac{1}{x+a}\)} & \((-1)^n \frac{n!}{(x+a)^{n+1}}\) &
符号交替，分子是 \(n!\) \\
\textbf{\((x+a)^\mu\)} & \(\mu(\mu-1)\cdots(\mu-n+1)(x+a)^{\mu-n}\) &
连乘 \(n\) 项 \\
\end{longtable}
}

\paragraph{二、微分}\label{ux4e8cux5faeux5206}

\begin{enumerate}
\def\labelenumi{\arabic{enumi}.}
\item
  定义：当自变量有个微小的变化 \(\Delta x\) 时，函数的真实变化量是
  \(\Delta y\)。如果 \(\Delta y\)
  可以表示为一个线性部分加上一个高阶无穷小，即
  \(\Delta y = A\Delta x + o(\Delta x)\)，那么这个线性部分 \(A\Delta x\)
  就叫作函数在这一点的微分，记作 \(dy=A\Delta x\)。
\end{enumerate}

\begin{itemize}
\item
  微分称为函数增量的\textbf{线性主部}
\end{itemize}

\begin{enumerate}
\def\labelenumi{\arabic{enumi}.}
\item
  可微和可导的关系：\textbf{可导 \(\iff\) 可微。}
  只要函数在某点有导数，它在该点就一定有微分；反之亦然。
\end{enumerate}

\paragraph{三、洛必达法则}\label{ux4e09ux6d1bux5fc5ux8fbeux6cd5ux5219}

\begin{enumerate}
\def\labelenumi{\arabic{enumi}.}
\item
  洛必达法则：
\end{enumerate}

每次使用前，必须严格确认以下三个条件，缺一不可：

\begin{enumerate}
\def\labelenumi{\arabic{enumi}.}
\item
  \textbf{未定式类型：} 在代入极限点 \(a\) 时，极限必须是标准的
  \(\frac{0}{0}\) 型或 \(\frac{\infty}{\infty}\) 型。
\item
  \textbf{去心邻域内可导：} 分子 \(f(x)\) 和分母 \(g(x)\) 在点 \(a\)
  的去心邻域内必须可导，并且分母的导数 \(g'(x) \neq 0\)。
\item
  \textbf{导数极限存在：} 求导后的极限
  \(\lim_{x \to a} \frac{f'(x)}{g'(x)}\)
  必须存在（或者是无穷大）。如果求导后的极限因为周期性震荡而``不存在''，则洛必达法则完全失效，不能由此反推原极限不存在。
\end{enumerate}

\textbf{核心公式：}

\[\lim_{x \to a} \frac{f(x)}{g(x)} = \lim_{x \to a} \frac{f'(x)}{g'(x)}\]

\paragraph{四、微分中值定理}\label{ux56dbux5faeux5206ux4e2dux503cux5b9aux7406}

\begin{enumerate}
\def\labelenumi{\arabic{enumi}.}
\item
  费马定理
\end{enumerate}

如果函数 \(f(x)\) 满足以下两个条件：

\begin{itemize}
\item
  \textbf{取得极值：} 在点 \(x_0\) 处取得极值（极大值或极小值）。
\item
  \textbf{存在导数：} 在点 \(x_0\) 处可导。
\end{itemize}

\textbf{结论：} 那么函数在该点的导数必定为零，即
\textbf{\(f'(x_0) = 0\)}。

\begin{enumerate}
\def\labelenumi{\arabic{enumi}.}
\item
  \textbf{罗尔定理}
\end{enumerate}

如果函数 \(f(x)\) 完美满足以下三个条件：

\begin{itemize}
\item
  \textbf{闭区间连续：} 在闭区间 \([a, b]\) 上连续。
\item
  \textbf{开区间可导：} 在开区间 \((a, b)\) 内可导。
\item
  \textbf{端点值相等：} \(f(a) = f(b)\)。
\end{itemize}

\textbf{结论：} 那么在开区间 \((a, b)\) 内，\textbf{至少存在一点
\(\xi\)}，使得 \textbf{\(f'(\xi) = 0\)}。

\begin{enumerate}
\def\labelenumi{\arabic{enumi}.}
\item
  \textbf{拉格朗日中值定理}
\end{enumerate}

如果函数 \(f(x)\) 满足以下两个条件（比罗尔定理少了一个端点相等的条件）：

\begin{itemize}
\item
  \textbf{闭区间连续：} 在闭区间 \([a, b]\) 上连续。
\item
  \textbf{开区间可导：} 在开区间 \((a, b)\) 内可导。
\end{itemize}

\textbf{结论：} 那么在开区间 \((a, b)\) 内，\textbf{至少存在一点
\(\xi\)}，使得以下等式成立：

\begin{itemize}
\item
  \textbf{形态一（斜率形态）：}

  \[f'(\xi) = \frac{f(b) - f(a)}{b - a}\]

  \emph{(几何意义：在曲线上至少有一点，其切线平行于连接两端点
  \(A(a, f(a))\) 和 \(B(b, f(b))\) 的割线。)}
\item
  \textbf{形态二（增量形态，考研最常用）：}

  \[f(b) - f(a) = f'(\xi)(b - a)\]

  \emph{(实战意义：把``函数值的差''瞬间转化成了``导数值乘上自变量的差''。这是降维打击的核心！)}
\end{itemize}

\begin{quote}
\textbf{洞察：} 可以把罗尔定理看作是拉格朗日中值定理的一个特例。当
\(f(a) = f(b)\) 时，割线是水平的，斜率为
0，拉格朗日的结论就自动变成了罗尔定理的 \(f'(\xi) = 0\)。
\end{quote}

\begin{enumerate}
\def\labelenumi{\arabic{enumi}.}
\item
  柯西中值定理
\end{enumerate}

如果有\textbf{两个}函数 \(f(x)\) 和 \(g(x)\)，它们同时满足以下条件：

\begin{itemize}
\item
  \textbf{闭区间连续：} 在闭区间 \([a, b]\) 上连续。
\item
  \textbf{开区间可导：} 在开区间 \((a, b)\) 内可导。
\item
  \textbf{分母导数不为零：} 对任意 \(x \in (a, b)\)，都有
  \(g'(x) \neq 0\)。（这是为了保证后面的分母有意义）
\end{itemize}

\textbf{结论：} 那么在开区间 \((a, b)\) 内，\textbf{至少存在一点
\(\xi\)}，使得以下等式成立：

\[\frac{f(b) - f(a)}{g(b) - g(a)} = \frac{f'(\xi)}{g'(\xi)}\]

\begin{quote}
\textbf{降维理解：} 如果您强行令 \(g(x) = x\)，那么
\(g(b) - g(a) = b - a\)，\(g'(\xi) = 1\)。此时柯西定理瞬间退化成了拉格朗日中值定理！这就证明了拉格朗日只是柯西的一个特例。
\end{quote}

\begin{enumerate}
\def\labelenumi{\arabic{enumi}.}
\item
  泰勒定理
\end{enumerate}

\begin{itemize}
\item
  \textbf{拉格朗日余项}：如果函数 \(f(x)\) 在含有 \(x_0\)
  的某个开区间内具有直到 \((n+1)\) 阶的导数，那么对于该区间内的任意一点
  \(x\)，都有：
\end{itemize}

\[f(x) = f(x_0) + f'(x_0)(x-x_0) + \frac{f''(x_0)}{2!}(x-x_0)^2 + \dots + \frac{f^{(n)}(x_0)}{n!}(x-x_0)^n + R_n(x)\]

其中，最关键的``尾巴'' \textbf{拉格朗日余项 \(R_n(x)\)} 为：

\[R_n(x) = \frac{f^{(n+1)}(\xi)}{(n+1)!}(x-x_0)^{n+1}\]

\emph{（这里的 \(\xi\) 严格介于 \(x_0\) 和 \(x\) 之间）}

\begin{quote}
\textbf{降维顿悟：} 假设 \(n=0\)，也就是只展开到第 0
阶。代入公式看看会发生什么？

\(f(x) = f(x_0) + \frac{f'(\xi)}{1!}(x-x_0)^1\)

移项一下：\(f(x) - f(x_0) = f'(\xi)(x-x_0)\)。

\textbf{这不就是拉格朗日中值定理吗！}
所以，泰勒定理就是把拉格朗日定理推广到了高阶导数的世界。
\end{quote}

\begin{itemize}
\item
  \textbf{皮亚诺型余项}：当我们在 \(x \to x_0\) 处对函数 \(f(x)\) 展开到
  \(n\) 阶时：
\end{itemize}

\[f(x) = f(x_0) + f'(x_0)(x-x_0) + \dots + \frac{f^{(n)}(x_0)}{n!}(x-x_0)^n + o((x-x_0)^n)\]

公式最后的这个 \textbf{\(o((x-x_0)^n)\)} 就是皮亚诺余项。

{\def\LTcaptype{none} % do not increment counter
\begin{longtable}[]{@{}llll@{}}
\toprule\noalign{}
\textbf{余项类型} & \textbf{形式} & \textbf{适用题型} &
\textbf{核心思维} \\
\midrule\noalign{}
\endhead
\bottomrule\noalign{}
\endlastfoot
\textbf{皮亚诺 (Peano)} & \(o(x^n)\) &
\textbf{求极限}（大题第一道、选择填空） &
局部逼近（极小范围），只看阶数，忽略细节 \\
\textbf{拉格朗日 (Lagrange)} & \(\frac{f^{(n+1)}(\xi)}{(n+1)!}x^{n+1}\)
& \textbf{证明题}（不等式放缩、中值定理压轴题） &
全局/区间逼近，必须精确量化误差大小 \\
\end{longtable}
}

\begin{itemize}
\item
  \textbf{麦克劳林公式}：将 \(x_0 = 0\)
  代入泰勒公式，即得到麦克劳林公式：

  \[f(x) = f(0) + f'(0)x + \frac{f''(0)}{2!}x^2 + \dots + \frac{f^{(n)}(0)}{n!}x^n + R_n(x)\]

  \begin{itemize}
  \item
    如果在做\textbf{极限计算}题，尾巴 \(R_n(x)\) 就写成皮亚诺余项
    \(o(x^n)\)。
  \item
    如果在做\textbf{中值定理证明}题，尾巴 \(R_n(x)\) 就写成拉格朗日余项
    \(\frac{f^{(n+1)}(\theta x)}{(n+1)!}x^{n+1}\) （其中
    \(0 < \theta < 1\)）
  \end{itemize}

  {\def\LTcaptype{none} % do not increment counter
  \begin{longtable}[]{@{}lll@{}}
  \toprule\noalign{}
  \textbf{函数 f(x)} & \textbf{展开式 (展开到常用的阶数)} &
  \textbf{记忆口诀与特征} \\
  \midrule\noalign{}
  \endhead
  \bottomrule\noalign{}
  \endlastfoot
  \textbf{\(e^x\)} &
  \(1 + x + \frac{1}{2!}x^2 + \frac{1}{3!}x^3 + o(x^3)\) &
  \textbf{全正号，有阶乘}。无穷级数的老祖宗。 \\
  \textbf{\(\sin x\)} &
  \(x - \frac{1}{3!}x^3 + \frac{1}{5!}x^5 + o(x^5)\) &
  奇函数，\textbf{只有奇数次幂，正负交替，有阶乘}。 \\
  \textbf{\(\cos x\)} &
  \(1 - \frac{1}{2!}x^2 + \frac{1}{4!}x^4 + o(x^4)\) &
  偶函数，\textbf{只有偶数次幂，正负交替，有阶乘}。 \\
  \textbf{\(\ln(1+x)\)} &
  \(x - \frac{1}{2}x^2 + \frac{1}{3}x^3 - \frac{1}{4}x^4 + o(x^4)\) &
  \textbf{正负交替，无阶乘！} (分母极易写错成阶乘) \\
  \textbf{\((1+x)^\alpha\)} &
  \(1 + \alpha x + \frac{\alpha(\alpha-1)}{2!}x^2 + o(x^2)\) &
  广义二项式展开，处理根号的绝对利器。 \\
  \textbf{\(\frac{1}{1-x}\)} & \(1 + x + x^2 + x^3 + o(x^3)\) &
  \textbf{全正号，无阶乘}。等比数列求和公式的逆用。 \\
  \textbf{\(\frac{1}{1+x}\)} & \(1 - x + x^2 - x^3 + o(x^3)\) & 把上式的
  \(x\) 换成 \(-x\) 即可，正负交替。 \\
  \textbf{\(\arctan x\)} &
  \(x - \frac{1}{3}x^3 + \frac{1}{5}x^5 + o(x^5)\) &
  奇函数，正负交替，\textbf{无阶乘}。 \\
  \end{longtable}
  }
\end{itemize}

\paragraph{五、单调性}\label{ux4e94ux5355ux8c03ux6027}

\begin{enumerate}
\def\labelenumi{\arabic{enumi}.}
\item
  单调性
\end{enumerate}

设函数 \(f(x)\) 在区间 \(I\) 上连续，在内部可导：

\begin{itemize}
\item
  \textbf{严格单调递增：} 如果在区间内 \(f'(x) > 0\)，则 \(f(x)\)
  单调递增（画出来是一直上坡）。
\item
  \textbf{严格单调递减：} 如果在区间内 \(f'(x) < 0\)，则 \(f(x)\)
  单调递减（画出来是一直下坡）。
\end{itemize}

\textbf{🚨 考场选择题极易踩坑的``命题陷阱''：}

\begin{itemize}
\item
  \textbf{陷阱 1：导数等于 0 会破坏严格单调性吗？}

  \begin{itemize}
  \item
    \textbf{结论：不会！} 如果 \(f'(x) \ge 0\)，且 \(f'(x) = 0\) （
    称为驻点）的点只是\textbf{离散的个别点}（比如 \(y = x^3\) 在 \(x=0\)
    处导数为
    0），那么函数依然是\textbf{严格单调递增}的。只有当导数在某一段``区间''内全都是
    0 时，函数才变成常数平移，失去严格单调性。
  \end{itemize}
\item
  \textbf{陷阱 2：单调区间的写法。}

  \begin{itemize}
  \item
    如果函数在相邻的两个区间都单调递增，且在分界点处连续，\textbf{单调区间可以合并}，并且端点通常写成闭区间（只要端点有定义且连续）。
  \end{itemize}
\end{itemize}

\paragraph{六、极值}\label{ux516dux6781ux503c}

\begin{enumerate}
\def\labelenumi{\arabic{enumi}.}
\item
  定义：设函数 \(f(x)\) 在点 \(x_0\) 的某个\textbf{邻域} \(U(x_0)\)
  内有定义（邻域就是以 \(x_0\) 为中心，向左向右延伸的一小段开区间
  \((x_0-\delta, x_0+\delta)\)）。

  \textbf{1. 极大值 (Local Maximum)}

  如果在该邻域 \(U(x_0)\) 内的\textbf{任意}一点
  \(x\)（\(x \neq x_0\)），都满足不等式：

  \[f(x) < f(x_0)\]

  那么就称 \(f(x_0)\) 是函数 \(f(x)\) 的一个\textbf{极大值}，点 \(x_0\)
  称为\textbf{极大值点}。

  \textbf{2. 极小值 (Local Minimum)}

  如果在该邻域 \(U(x_0)\) 内的\textbf{任意}一点
  \(x\)（\(x \neq x_0\)），都满足不等式：

  \[f(x) > f(x_0)\]

  那么就称 \(f(x_0)\) 是函数 \(f(x)\) 的一个\textbf{极小值}，点 \(x_0\)
  称为\textbf{极小值点}。

  \emph{(注：上述定义是``严格极值''。如果把不等号换成 \(\le\) 或
  \(\ge\)，就是一般极值。在考研中，默认讨论的都是严格极值。)}
\item
  \textbf{必要条件}：如果函数 \(f(x)\) 在 \(x_0\)
  处取得极值，\textbf{并且 \(f(x)\) 在 \(x_0\) 处可导}，那么必定有
  \textbf{\(f'(x_0) = 0\)}。
\item
  \textbf{充分条件}
\end{enumerate}

\begin{itemize}
\item
  \textbf{第一充分条件：看一阶导数的``左右变号''}

  设函数\(f(x)\)在\(x_0\)的某领域\((x_0-\delta,x_0+\delta)\)内\textbf{连续}，\(x_0\)是驻点（\(f'(x_0) = 0\)）或不可导点

  \textbf{判定方法：} 观察点 \(x_0\) 左右两侧一阶导数 \(f'(x)\) 的符号。

  \begin{itemize}
  \item
    \textbf{左正右负 (\(+ \to -\)):} 曲线在 \(x_0\) 左边上坡，右边下坡
    \(\implies\) \textbf{极大值点}。
  \item
    \textbf{左负右正 (\(- \to +\)):} 曲线在 \(x_0\) 左边下坡，右边上坡
    \(\implies\) \textbf{极小值点}。
  \item
    \textbf{符号不变 (\(+ \to +\) 或 \(- \to -\)):} 一直上坡或一直下坡
    \(\implies\) \textbf{不是极值点}（只是个过路的停顿点）。
  \end{itemize}
\item
  \textbf{第二充分条件：看二阶导数的``正负''}

  \textbf{判定方法：} 已知 \(f'(x_0) = 0\)且二阶导数存在，将 \(x_0\)
  代入二阶导数 \(f''(x)\) 中：

  \begin{itemize}
  \item
    \textbf{\(f''(x_0) < 0\)：} 二阶导小于
    0，曲线在该点是``凸''的，此时处于最高处 \(\implies\)
    \textbf{极大值点}。
  \item
    \textbf{\(f''(x_0) > 0\)：} 二阶导大于
    0，曲线在该点是``凹''的，此时处于最低处 \(\implies\)
    \textbf{极小值点}。
  \end{itemize}
\end{itemize}

\paragraph{七、凹凸性}\label{ux4e03ux51f9ux51f8ux6027}

\begin{enumerate}
\def\labelenumi{\arabic{enumi}.}
\item
  凹凸性的定义：设函数\(f(x)\)在区间\(I\)上连续，如果在区间内任意取两点
  \(x_1, x_2\)，恒有：

  \begin{itemize}
  \item
    \(f\left(\frac{x_1+x_2}{2}\right) < \frac{f(x_1)+f(x_2)}{2} \implies\)
    \textbf{中点的值低于两端点的平均值} \(\implies\) 曲线向下凹
  \item
    \(f\left(\frac{x_1+x_2}{2}\right) > \frac{f(x_1)+f(x_2)}{2} \implies\)
    \textbf{中点的值高于两端点的平均值} \(\implies\) 曲线向上凸
  \end{itemize}
\item
  凹凸性的判别：若函数\(f(x)\)在区间\(I\)上\(f''(x)<0(>0)\)，则曲线\(y=f(x)\)在\(I\)上凹（凸）的。
\item
  拐点：
\end{enumerate}

\begin{itemize}
\item
  定义：设曲线 \(y = f(x)\) 在点 \((x_0, f(x_0))\)
  处连续。如果曲线在该点左右两侧的\textbf{凹凸性相反}，那么点
  \((x_0, f(x_0))\) 就称为曲线的拐点。
\item
  必要条件：在函数的定义域内，点\((x_0,f(x_0))\)是拐点的必要条件是\(f''(x) = 0\)
  或者\(f''(x)\) 不存在
\item
  充分条件：如果
  \textbf{\(f''(x_0) = 0\)}且\(f'''(x_0) \neq 0\)，并且该点领域内三阶导数存在，那么
  点 \((x_0, f(x_0))\)是曲线\(y=f(x)\)的拐点
\end{itemize}

\paragraph{八、渐近线}\label{ux516bux6e10ux8fd1ux7ebf}

\begin{enumerate}
\def\labelenumi{\arabic{enumi}.}
\item
  定义：当曲线上的动点沿着曲线无限远离原点时，若动点与某条直线的距离趋近于0，则称该直线为曲线的渐近线。
\item
  水平渐近线：若\(\lim_{x \to \infty} f(x)=C\),=，那么直线
  \textbf{\(y = C\)} 是曲线\(y=f(x)\)的一条水平渐近线。
\item
  铅直渐近线：若
  \(\lim_{x \to x_0} f(x)=\infty\)（有时需要分左右极限），则\(x=x_0\)是曲线的一条铅直渐近线。
\item
  斜渐近线：设函数 \(y = f(x)\) 定义在无限延伸的区间上。如果存在两个常数
  \textbf{\(a\)} 和 \textbf{\(b\)}（其中 \textbf{\(a \neq 0\)}），使得当
  \(x \to +\infty\)（或 \(x \to -\infty\)）时，函数 \(f(x)\) 与直线
  \(y = ax + b\) 之间的\textbf{差值趋近于 0}，即：

  \[\lim_{x \to \infty} [f(x) - (ax + b)] = 0\]

  那么，我们就称直线 \textbf{\(y = ax + b\)} 是曲线 \(y = f(x)\)
  的一条\textbf{斜渐近线}。
\end{enumerate}

\paragraph{九、曲率}\label{ux4e5dux66f2ux7387}

\begin{enumerate}
\def\labelenumi{\arabic{enumi}.}
\item
  定义：设曲线 \(C\) 是一条光滑曲线。我们在曲线上取一点
  \(M\)，并以此为基准点。
\end{enumerate}

再在曲线上取邻近的一点 \(M'\)。

\begin{itemize}
\item
  设从 \(M\) 到 \(M'\) 的一段\textbf{弧长}为
  \textbf{\(\Delta s\)}（您在轨道上实际开过的距离）。
\item
  设在这两点处，曲线的\textbf{切线转过的角度}为
  \textbf{\(\Delta \alpha\)}（您的方向盘实际打过的角度）。
\end{itemize}

\textbf{定义：} 当点 \(M'\) 沿着曲线无限逼近点 \(M\)（即
\(\Delta s \to 0\)）时，切线转角的绝对值 \(|\Delta \alpha|\) 与弧长
\(\Delta s\) 之比的极限，就定义为曲线在点 \(M\) 处的\textbf{曲率
(Curvature)}，记作 \(K\)。用严谨的微积分符号表示就是：

\[K = \lim_{\Delta s \to 0} \left| \frac{\Delta \alpha}{\Delta s} \right| = \left| \frac{d\alpha}{ds} \right|\]

\begin{enumerate}
\def\labelenumi{\arabic{enumi}.}
\item
  计算公式：已知
  \(K = \left| \frac{d\alpha}{ds} \right|\)。我们需要把它转化成只和
  \(x, y\) 有关的式子。

  \textbf{1. 搞定分子 \(d\alpha\)（切线倾角的微分）：}

  切线的斜率就是导数，即 \(\tan \alpha = y'\)。

  两边取反三角函数：\(\alpha = \arctan y'\)。

  对 \(x\) 求导（复合函数求导）：

  \[\frac{d\alpha}{dx} = \frac{1}{1 + (y')^2} \cdot y''\]

  所以，\textbf{\(d\alpha = \frac{y''}{1 + (y')^2} dx\)}。

  \textbf{2. 搞定分母 \(ds\)（弧长微分）：}

  根据勾股定理，极小的一段弧长 \(ds\) 可以看作横向微元 \(dx\) 和纵向微元
  \(dy\) 构成的直角三角形的斜边：

  \[ds = \sqrt{(dx)^2 + (dy)^2} = \sqrt{1 + \left(\frac{dy}{dx}\right)^2} dx = \sqrt{1 + (y')^2} dx\]

  \textbf{3. 上下相除}

  把算出来的 \(d\alpha\) 和 \(ds\) 代入曲率定义式
  \(K = \left| \frac{d\alpha}{ds} \right|\) 中：

  \[K = \frac{\left| \frac{y''}{1 + (y')^2} dx \right|}{\sqrt{1 + (y')^2} dx} = \frac{|y''|}{(1 + (y')^2) \cdot \sqrt{1 + (y')^2}}\]

  合并分母，瞬间得到我们死记硬背的考场公式：

  \[K = \frac{|y''|}{(1 + (y')^2)^{\frac{3}{2}}}\]
\item
  曲率半径：曲率半径（通常记作 \(R\) 或 \(\rho\)）就是\textbf{曲率 \(K\)
  的倒数}。

  \[R = \frac{1}{K}\]
\end{enumerate}

{\def\LTcaptype{none} % do not increment counter
\begin{longtable}[]{@{}
  >{\raggedright\arraybackslash}p{(\linewidth - 4\tabcolsep) * \real{0.3333}}
  >{\raggedright\arraybackslash}p{(\linewidth - 4\tabcolsep) * \real{0.3333}}
  >{\raggedright\arraybackslash}p{(\linewidth - 4\tabcolsep) * \real{0.3333}}@{}}
\toprule\noalign{}
\textbf{对比维度} & \textbf{⛰️ 极值点 (一阶特征)} & \textbf{〰️ 拐点
(二阶特征)} \\
\midrule\noalign{}
\endhead
\bottomrule\noalign{}
\endlastfoot
\textbf{直观几何意义} & 曲线局部位置的 \textbf{最高处} 或
\textbf{最低处}。 \emph{(上坡变下坡，或下坡变上坡)} & 曲线弯曲姿态的
\textbf{分水岭}。 \emph{(碗口向上变向下，或向下变向上)} \\
答题最终形式 & 是一个横坐标 \textbf{数值}：\(x = x_0\) & 是一个二维
\textbf{坐标点}：\((x_0, f(x_0))\) \\
必要条件 & 令 \textbf{一阶导数 \(f'(x) = 0\)} 的点 (驻点) 或
\textbf{一阶导数不存在} 的点。 & 令 \textbf{二阶导数 \(f''(x) = 0\)}
的点 或 \textbf{二阶导数不存在} 的点。 \\
充分条件一 & \vtop{\hbox{\strut 看 \textbf{一阶导数 \(f'(x)\)}
左右是否变号。}\hbox{\strut • 左 \(+\) 右 \(-\)：极大值点
}\hbox{\strut • 左 \(-\) 右 \(+\)：极小值点}\hbox{\strut  •
不变号：不是极值点}} & \vtop{\hbox{\strut 看 \textbf{二阶导数
\(f''(x)\)} 左右是否变号。 }\hbox{\strut • 只要变号
(不管从正变负还是负变正)：\textbf{就是拐点} }\hbox{\strut •
不变号：不是拐点}} \\
充分条件二 & \vtop{\hbox{\strut 针对一阶驻点 \(f'(x_0)=0\)： 看
\textbf{二阶导数 \(f''(x_0)\)} 的符号。 }\hbox{\strut •
\(< 0\)：极大值点 (\(\frown\)) }\hbox{\strut • \(> 0\)：极小值点
(\(\smile\))}} & \vtop{\hbox{\strut 针对二阶驻点 \(f''(x_0)=0\)： 看
\textbf{三阶导数 \(f'''(x_0)\)}。}\hbox{\strut  • 只要
\textbf{\(\neq 0\)}：\textbf{就是拐点}}} \\
充分条件三 & 若前 \(n-1\) 阶导数均为 0， 且第 \(n\) 阶导数
\(f^{(n)}(x_0) \neq 0\)： 当 \(n\) 为 \textbf{偶数} 时，是极值点。 &
若前 \(n-1\) 阶导数均为 0， 且第 \(n\) 阶导数 \(f^{(n)}(x_0) \neq 0\)：
当 \(n\) 为 \textbf{奇数} 时，是拐点。 \\
特例 & \(y = x^3\) 在 \(x=0\) 处： \(y'=0\)，但左右一阶导不变号，
\textbf{不是极值点}。 & \(y = x^4\) 在 \(x=0\) 处：
\(y''=0\)，但左右二阶导不变号， \textbf{不是拐点}。 \\
\textbf{闭区间特例} & 区间端点 (\(a\) 或 \(b\)) \textbf{绝对不能}
称为极值点。(极值必须有双侧邻域) & 区间端点同样 \textbf{不能}
称为拐点。(拐点也要求左右两侧凹凸性不同) \\
\end{longtable}
}

\end{document}
